\section{改动概述}
\label{sec:modifications}

在开发 DreamerV2 时，我们以 Dreamer 智能体\citep{hafner2019dreamer}为起点。本小节概述了为在 Atari 基准上取得高性能而引入的改动，也列出了曾经尝试、但未带来明显提升，因此没有纳入 DreamerV2 的设计。

以下改动在实验中被证明有帮助：

\begin{itemize}
\itempar{分类潜变量} 在世界模型中使用分类潜状态，并通过直通梯度进行训练；相比之下，原始 Dreamer 使用的是带重参数化梯度的高斯潜变量。
\itempar{KL 平衡} 在 KL 损失中分别缩放先验交叉熵和后验熵，以鼓励模型学习准确的时序先验，而不是使用 free nats。
\itempar{仅使用 Reinforce} 对 Atari 而言，Reinforce 梯度明显优于通过动力学反向传播得到的梯度；而在连续控制任务中，动力学反向传播的效果要好得多。
\itempar{模型规模} 增加所有模型组件中每层单元或特征图的数量，使参数量从 13M 增至 22M。
\itempar{策略熵} 在想象训练和数据收集过程中都用策略熵正则化来促进探索，而不是只在数据收集时加入外部动作噪声。
\end{itemize}

以下改动也被尝试过，但没有带来实质性收益：

\begin{itemize}
\itempar{二值潜变量} 用更多二值潜变量替代分类潜变量，本来可能促进更解耦的表征，但实际效果更差。
\itempar{长期熵} 将策略熵纳入价值函数的时序差分损失，使 actor 在规划视野之外也倾向于寻找动作熵较高的状态。
\itempar{混合 actor 梯度} 将 Reinforce 梯度与动力学反向传播梯度结合来训练 actor，相比只使用 Reinforce 仅带来很小收益，或没有收益。
\itempar{调度策略} 对学习率、KL 系数、actor 熵损失系数以及 actor 梯度混合比例（从 0.1 降至 0）进行调度，只带来边际收益或没有收益。
\itempar{层归一化} 在 RSSM 潜在转移模型中的 GRU 内使用层归一化，而不是不做归一化，收益很小或没有收益。
\end{itemize}

由于计算开销很大，我们无法对上述所有改动进行完整的消融研究。完整评估需要 55 个任务、5 个随机种子，并且每项改动运行 10 天，折合每项改动超过 60,000 个 GPU 小时。因此，正文只纳入了最重要设计选择的消融结果。
